\documentclass[10pt,a4paper]{amsart}
\usepackage[margin=19mm]{geometry}
\usepackage{amsmath,amssymb,mathtools}
\usepackage[scr=rsfs,cal=euler]{mathalfa}
\usepackage{xfrac}
\usepackage[unicode,hidelinks,bookmarks=false]{hyperref}
\newcommand{\ie}{i.e.\ }
\newcommand{\cf}{cf.\ }
\newcommand{\resp}{respectively\ }
\newcommand{\Subsection}{\subsection}
\providecommand{\nobreakdash}{\nobreak\hbox{-}}
\setlength{\emergencystretch}{2em}
\multlinegap0pt
\usepackage{bm}
\usepackage{bbm}
\usepackage{dsfont}
\usepackage{xspace}
\usepackage{cancel}
\usepackage{mathtools}
\usepackage{cleveref}
\usepackage{amsthm}
\makeatletter
\@ifundefined{thm}{\newtheorem{thm}{Thm}}{}
\@ifundefined{lem}{\newtheorem{lem}[thm]{Lemma}}{}
\makeatother
\makeatletter
\newcommand{\C}{\mathbb{C}}
\expandafter\ifx\csname lst\endcsname\relax
\newenvironment{lst}{%
\refstepcounter{equation}%
\begin{center}
\begin{minipage}{.9\textwidth}}{%
\end{minipage}%
\makebox[.1\textwidth][r]{(\theequation)}%
\end{center}
}
\fi
\newcommand\restr[2]{{
	\left.\kern-\nulldelimiterspace
	#1
	\vphantom{\big|}
	\right|_{#2}
	}}
\makeatother
\title[Algebraic Hodge generic points are dense]{Algebraic Hodge generic points are dense}
\author{Gregorio Baldi, Gal Binyamini,, David Urbanik}
\date{Source: arXiv:2606.08882v2 (CC BY 4.0)}
\begin{document}
\maketitle

\noindent\emph{Source and licence.} Algebraic Hodge generic points are dense. Version arXiv:2606.08882v2 (CC BY 4.0), distributed under the Creative Commons Attribution 4.0 International licence, as recorded in the arXiv licence metadata for this exact version and shown on the abstract page https://arxiv.org/abs/2606.08882v2. Full rights notice: licenses/arxiv-2606.08882-CC-BY-4.0.txt.\par\smallskip

\noindent\textit{Editorial scope.} Counted unit: the Lem whose source label is \texttt{deeplemma} in the article (source0.tex lines 1140--1152, source label \texttt{deeplemma}), delivered together with the article's own complete original proof of it (source0.tex lines 1154--1170, one \texttt{proof} environment of 17 lines). No mathematics is written, repaired or re-derived: statement and proof are the article's own text line for line. Required context retained uncounted: BGgrapheq (lines 1081--1084); lemmacodim (lines 1174--1176). Every definition, display and label of those lines is retained unchanged. Omitted from this package and not redistributed: 1--1080, 1085--1139, 1153--1153, 1171--1173, 1177--1588. Nothing inside a retained excerpt is omitted or altered: every retained body is delivered exactly as the article's lines are written, so no omission receipt of any kind accompanies this package. Citations: the retained excerpts carry no citation command at all, so no bibliography entry of the article is transcribed and no reference is redistributed. Reference adaptation: none was needed and none was made; every reference inside the delivered unit resolves inside it, so no unresolved reference or citation is produced. Printed slips kept literal: no printed word, symbol, hypothesis or number of the counted statement or of its proof was altered. Layout and class: the article's own class is \texttt{article} with packages including amsmath, amsfonts, amssymb, framed, mathtools, enumerate, mathrsfs, hyperref, enumitem, scrextend, amsthm, bbold, bbm, comment; the wrapper is the lane's pinned \texttt{amsart} editorial wrapper at 10pt on A4, which restates the theorem-like environments the retained text needs and adds the article's own macro definitions that the retained text needs (\texttt{\textbackslash C}, \texttt{\textbackslash restr}), with extra wrapper packages (bm, bbm, dsfont, xspace, cancel, mathtools, cleveref). The delivered numbering is the unit's own. Assets: no retained span contains a figure environment, a graphics-inclusion command, a PDF-page inclusion, a table or a TikZ picture; no image file is copied into this package, the package declares no asset dependency and no image file is redistributed with it. Licence: version arXiv:2606.08882v2 of this article is distributed under the Creative Commons Attribution 4.0 International licence, as recorded in the arXiv licence metadata for this exact version and shown on the abstract page https://arxiv.org/abs/2606.08882v2. A full rights notice is delivered at licenses/arxiv-2606.08882-CC-BY-4.0.txt. Characters: every retained excerpt is pure ASCII, so the delivered engine, XeLaTeX with its default Latin Modern text font, reports no missing character.
\begin{equation}
\label{BGgrapheq}
s \mapsto \{ (\mathbf{J}_{j}(s)) \}^{n^2}_{j=0} 
\end{equation}

\begin{lem}\label{lemmacodim}
    Let $(A,\mathfrak{m})$ be a local ring (which is a $\C$-algebra) with spectrum $B$, and $C$ be the formal spectrum of an irreducible formal power series ring over $\C$. Let $R$ be a non-zero element of the coordinate ring of $B\times C$ such that $\restr{R}{\mathfrak{m}\times C}$ is not identically zero. If $B$ has pure dimension, then $H=V(R)$ has codimension one in $B\times C$. 
\end{lem}

\begin{lem}
\label{deeplemma}
For each $j = 0, \hdots, n^2$, fix a closed embedding $B_{\mathbf{J}}(j) \subset \mathbb{P}^1 \times \mathbb{A}^{r}$ over $\mathbb{P}^1$, and let $\mathcal{B}_{\mathbf{J}} \subset \mathbb{P}^1 \times \mathbb{A}^{Nr}$ be the associated embedding of $X := \mathcal{B}_{\mathbf{J}}$. Let $Y \subset \mathbb{A}^{n^2}$ be an irreducible affine variety defined over $K$, and write $\pi_{j} : \mathcal{B}_{\mathbf{J}} \to B_{\mathbf{J}}(j)$ for the obvious projections.


Suppose we are given: 
\begin{itemize}
    \item a point $s \in \mathbb{P}^1(K)$;
    \item a point $E = (E_{1}, E_{2}) \in (X_{s} \times Y)(\mathbb{C})$ with $\overline{E_{2}}^{K-\operatorname{Zar}} = Y$; and
    \item $K$-algebraic relations $R_{0}, \hdots, R_{n^2}$, with each $R_{j}$ a relation on $B_{\mathbf{J}}(j)_{s} \times Y$, such that $R_{j}$ vanishes at $(\pi_{j}(E_{1}), E_{2})$ for each $j$ and is not zero in the stalk of $B_{\mathbf{J}}(j)_{s} \times Y$ at $(\pi_{j}(E_{1}), E_{2})$.
\end{itemize}
Then the projection to $X_{s}$ of the vanishing locus $V(\pi_{0}^{*} R_{0}, \hdots, \pi^{*}_{n^2} R_{n^2})$ contains a component of codimension at least $N - \dim Y$ in $X_{s}$, containing the image of $E$, and defined by polynomials whose degrees under the fixed affine embedding are bounded by a polynomial (depending only on the fixed embeddings, and in particular independent of $s$) in the degrees of the $R_{0}, \hdots, R_{n^2}$.
\end{lem}

\begin{proof}
From now on we write $R_{i} = \pi^{*}_{i} R_{i}$. We claim that the locus $V(R_{0}, \hdots, R_{n^2})$ has a component containing $E = (E_{1}, E_{2})$ and which has codimension $N$. As explained in the paragraph containing \eqref{BGgrapheq}, $\mathcal{B}_{\mathbf{J}}$ is a torsor, so it follows that $\mathcal{B}_{\mathbf{J},s}$ is a smooth variety, and hence locally irreducible. To compute the dimension of $$V(R_{0}, \hdots, R_{n^2})$$ at $E$ one can work inside the ring of formal germs of $X_{s} \times Y$ at $E$, which is naturally the quotient of a ring of the form
\begin{displaymath}
  M:=   \mathbb{C}[[\{ v^{i}_{0} \}^{r}_{i=1}, \hdots, \{ v^{i}_{n^2} \}^{r}_{i=1}, y_{1}, \hdots, y_{\mu} ]]
\end{displaymath}
where the $\{ v^{i}_{j} \}^{r}_{i=1}$ are coordinates at $\pi_{i}(E)$ on the affine space $\mathbb{A}^{r}$ containing $\mathcal{B}_{\mathbf{J}}(j)_{s}$, and the $y_{1}, \hdots, y_{\mu}$ are formal coordinates for $Y$ at $E_{2}$. In this ring, each $R_{j}$ takes the form $R_{j}(\{ v^{i}_{j} \}^{r}_{i=1}, y_{1}, \hdots, y_{\mu})$, and we claim that $R_j$ does not lie in the ideal $(y_1,\dots, y_\mu) \subset M$. Indeed, if this were true, it would already be true in the global affine coordinate ring of $X_s\times Y$, but, since the relations $R_j$ are defined over $K$, this would contradict the fact that the $K$-Zariski closure of $E_{2}$ is $Y$.



Recall that each $R_j$ is a $K$-algebraic relation on $B_{\mathbf{J}}(j)_{s} \times Y$ vanishing at $(\pi_j(E_1),E_2)$ and not zero in the stalk of $B_{\mathbf{J}}(j)_{s} \times Y$ at $(\pi_{j}(E_{1}), E_{2})$.  It then follows by iteratively applying \Cref{lemmacodim} below that $(R_{0}, \hdots, R_{n^2})$ is a regular sequence. Indeed, we can argue by induction as follows. The base case is given by considering the relation $R_0$ on $\mathcal{B}_{\mathbf{J}, {s}}\times B_Q$ which is not identically zero, and so the associated $V(R_0)$ has codimension one in a neighbourhood of $(E_1, E_2)$. Let $C$ be the formal stalk at $\pi_j(E_1)$ of $B_{\mathbf{J}}(j)_{s}$. Let $B$ be the formal stalk at $\mathfrak{m}=(\pi_i(E_1)_{i\neq j},E_2)$ of 
\begin{displaymath}
V ((R_i)_{i< j} ) \subset  B_Q \times \prod_{i\neq j} B_{\mathbf{J}}(i)_{s} .
\end{displaymath}
 By the inductive assumption, $B$ is of pure dimension and  $R_j$ is not identically zero on $\{ \mathfrak{m} \} \times C$. \Cref{lemmacodim} then shows that $V(R)$ has codimension one in $B\times C$. Since $(R_{0}, \hdots, R_{n^2})$ is a regular sequence, the ideal cut out by these functions defines a locus of codimension at least $N$ in some neighborhood of $(E_1,E_2)$ in $X_s\times Y$. It follows that its projection to $X_{s}$ has local codimension at least $N - \dim Y$. 

The claim on the degrees is an easy consequence of elimination theory.
\end{proof}

\end{document}
